% GENERATED FILE. Edit canonical lessons, then run npm run generate:books.
% canonical-source-hash: fnv1a64:3a0b6af9c8ece551
% canonical-lessons: ZH-C516-luoye, ZH-C516-shu-tree, ZH-W516-shu-tree, ZH-C516-guoshu, ZH-C516-shulin, ZH-C516-chengshu

\chapter{Trees and Leaves}
\label{ch:zh-trees-and-leaves}
\hlchaptermodality{516}

\begin{chapteropening}
\textbf{By the end of this chapter:} \emph{I can talk about trees, fruit trees, woods, fallen leaves and fruit that is not ripe yet.}

Complete the last lesson of chapter 516: I can talk about trees, fruit trees, woods, fallen leaves and fruit that is not ripe yet.
\end{chapteropening}

\section[luòyè]{luòyè — fallen leaves}

\label{lesson:ZH-C516-luoye}



\begin{quote}
Before the new one: say the Mandarin for to escape the summer heat, then the Mandarin for to bring in the harvest.
\end{quote}

\subsection*{You'll want to know: luòyè}
\emph{luòyè} — \textquotedblleft{}fallen leaves\textquotedblright{}.

\begin{quote}
\emph{qiūtiān gōngyuán lǐ dōu shì luòyè} — In autumn the park is full of fallen leaves.
\end{quote}

Its characters come later; for now it is a word for the ear.

\begin{cousinweb}
\emph{luòyè} is built from \emph{luò} 'to fall' and \emph{yè} 'a leaf', both inherited from Middle Chinese, \emph{lak} and \emph{yep} (Guangyun, in Baxter's transcription). The dictionary says the first is the word for what trees do when their leaves come down, and glosses the second 'branches and leaves'.
\end{cousinweb}

\begin{grammarlens}[title={What you've built}]
One more word for saying how it was.
\end{grammarlens}

\subsection*{Your turn}
\begin{itemize}
\raggedright
  \item \emph{Say it:} \emph{luòyè}
  \item \emph{Say it:} \emph{luòyè}, once more
  \item \emph{Say it:} \emph{qiūtiān gōngyuán lǐ dōu shì luòyè} — In autumn the park is full of fallen leaves.
  \item \emph{Recall:} say the Mandarin for to escape the summer heat, then the Mandarin for to bring in the harvest, then say \emph{luòyè} again
\end{itemize}

\begin{tcolorbox}[breakable,colback=teal!4,colframe=teal!35!black,title={Before you move on}]
What does luòyè mean? (\textquotedblleft{}Fallen leaves\textquotedblright{}.) Say it once more.
\end{tcolorbox}
\section[\texorpdfstring{\zh{树} (shù)}{树 (shù)}]{\zh{树} (shù) — tree}

\label{lesson:ZH-C516-shu-tree}



\begin{quote}
Before the new one: say the Mandarin for to bring in the harvest, then the Mandarin for fallen leaves.
\end{quote}

\subsection*{You'll want to know: \zh{树}}
\textbf{\zh{树}} — \emph{shù} — \textquotedblleft{}tree\textquotedblright{}.

\begin{quote}
\textbf{\zh{秋天的树最好看}} \emph{qiūtiān de shù zuì hǎokàn} — Trees look their best in autumn.
\end{quote}

The new character, in pieces: \zh{树} — nine strokes in three pieces: on the left the tree piece, four strokes, as in \zh{机}; in the middle \zh{又}, two strokes, the second shrunk to a short stroke; and on the right a bar, a vertical that hooks to the left, and a dot. The next lesson puts it in your hand.

\begin{cousinweb}
\textbf{\zh{树}} (\emph{shù}) is inherited from Middle Chinese \emph{dzyuH} (Guangyun, in Baxter's transcription), glossed as the general name for trees and 'to set up'; the dictionary also lists \emph{dzyuX}, glossed 'to prop up'. The falling tone of shù comes from \emph{dzyuH}.
\end{cousinweb}

\begin{grammarlens}[title={What you've built}]
One more word for saying how it was.
\end{grammarlens}

\subsection*{Your turn}
\begin{itemize}
\raggedright
  \item \emph{Say it:} \emph{shù}
  \item \emph{Say it:} \emph{shù}, once more
  \item \emph{Say it:} \emph{qiūtiān de shù zuì hǎokàn} — Trees look their best in autumn.
  \item \emph{Recall:} say the Mandarin for to bring in the harvest, then the Mandarin for fallen leaves, then say \emph{shù} again
\end{itemize}

\begin{tcolorbox}[breakable,colback=teal!4,colframe=teal!35!black,title={Before you move on}]
What does \zh{树} mean? (\textquotedblleft{}Tree\textquotedblright{}.) Say it once more.
\end{tcolorbox}
\section[\texorpdfstring{\zh{树} (shù)}{树 (shù)}]{\zh{树} — the tree piece, \zh{又}, and three strokes on the right}

\label{lesson:ZH-W516-shu-tree}



\begin{quote}
The word first, then the character inside it.

\begin{itemize}
\raggedright
  \item \emph{Recall:} say \emph{shù} — \textbf{R1}, one lesson back, before it decays
\end{itemize}

\begin{quote}
\zh{机} \zh{又}
\end{quote}

\emph{Read it:} each one aloud, with nothing said first — \emph{jī}, \emph{yòu}
\end{quote}

\subsection*{Script — the tree piece, \zh{又}, and three strokes on the right}
\hlblockfigure{\detokenize{figures/ZH-W516-shu-tree-filmstrip.pdf}}{How \zh{树} is written, stroke by stroke}

\begin{quote}
\zh{树}
\end{quote}

\textbf{Three pieces, left to right.} Count the pieces, not the strokes.

\textbf{The left piece, the tree piece, four strokes}, narrow, as in \zh{机}: a short bar, a vertical down through it, a stroke falling to the left from the crossing, then a short dot to the right.

\textbf{The middle piece, \zh{又}, two strokes}: a short bar that turns into a stroke falling to the left, then a short stroke falling to the right across it.

\textbf{The right piece, three strokes}: a bar, a long vertical down the right side that hooks to the left at the foot, then a dot on the left of the vertical.

4 + 2 + 3 = nine.

\subsection*{Your turn}
\begin{itemize}
\raggedright
  \item \emph{Write it:} \zh{树}, all nine, with the model in front of you
  \item \emph{Cover:} write it again from memory, and count the pieces, not the strokes
\end{itemize}

\begin{tcolorbox}[breakable,colback=teal!4,colframe=teal!35!black,title={Before you move on}]
How many pieces does \zh{树} have? (\textbf{Three: the tree piece, \zh{又} and a right piece}.) What is the last stroke? (\textbf{A dot beside the hooked vertical}.)

Source: \href{https://raw.githubusercontent.com/chanind/hanzi-writer-data/68d10a4b21150cae5e1ebbd223eed289cf32d90c/data/\%E6\%A0\%91.json}{Hanzi Writer Data \zh{树}, pinned stroke paths and medians}.
\end{tcolorbox}
\section[\texorpdfstring{\zh{果树} (guǒshù)}{果树 (guǒshù)}]{\zh{果树} (guǒshù) — fruit tree}

\label{lesson:ZH-C516-guoshu}



\begin{quote}
Before the new one: say the Mandarin for fallen leaves, then the Mandarin for a tree.
\end{quote}

\subsection*{You'll want to know: \zh{果树}}
\textbf{\zh{果树}} — \emph{guǒshù} — \textquotedblleft{}fruit tree\textquotedblright{}.

\begin{quote}
\textbf{\zh{我外公家有一些果树}} \emph{wǒ wàigōng jiā yǒu yìxiē guǒshù} — There are some fruit trees at my grandad's.
\end{quote}

\begin{cousinweb}
\textbf{\zh{果树}} (\emph{guǒshù}) is built from \zh{果} 'fruit' and \zh{树} 'a tree', both inherited from Middle Chinese, \emph{kwaX} and \emph{dzyuH} (Guangyun, in Baxter's transcription), which glosses the first 'bold' and 'firm', and also 'the fruit of a tree'.
\end{cousinweb}

\begin{grammarlens}[title={What you've built}]
One more word for saying how it was.
\end{grammarlens}

\subsection*{Your turn}
\begin{itemize}
\raggedright
  \item \emph{Say it:} \emph{guǒshù}
  \item \emph{Say it:} \emph{guǒshù}, once more
  \item \emph{Say it:} \emph{wǒ wàigōng jiā yǒu yìxiē guǒshù} — There are some fruit trees at my grandad's.
  \item \emph{Recall:} say the Mandarin for fallen leaves, then the Mandarin for a tree, then say \emph{guǒshù} again
\end{itemize}

\begin{tcolorbox}[breakable,colback=teal!4,colframe=teal!35!black,title={Before you move on}]
What does \zh{果树} mean? (\textquotedblleft{}Fruit tree\textquotedblright{}.) Say it once more.
\end{tcolorbox}
\section[shùlín]{shùlín — woods; a grove}

\label{lesson:ZH-C516-shulin}



\begin{quote}
Before the new one: say the Mandarin for a tree, then the Mandarin for a fruit tree.
\end{quote}

\subsection*{You'll want to know: shùlín}
\emph{shùlín} — \textquotedblleft{}woods; a grove\textquotedblright{}.

\begin{quote}
\emph{wǒ hé sūnnǚ zài shùlín lǐ zǒu le liǎng ge xiǎoshí} — My granddaughter and I walked in the woods for two hours.
\end{quote}

Its characters come later; for now it is a word for the ear.

\begin{cousinweb}
\emph{shùlín} is built from \zh{树} 'a tree' and \emph{lín} 'a wood', both inherited from Middle Chinese, \emph{dzyuH} and \emph{lim} (Guangyun, in Baxter's transcription), which says, after an older dictionary, that land out in the open is called \emph{lim}, and after another, that flat land with stands of trees is.
\end{cousinweb}

\begin{grammarlens}[title={What you've built}]
One more word for saying how it was.
\end{grammarlens}

\subsection*{Your turn}
\begin{itemize}
\raggedright
  \item \emph{Say it:} \emph{shùlín}
  \item \emph{Say it:} \emph{shùlín}, once more
  \item \emph{Say it:} \emph{wǒ hé sūnnǚ zài shùlín lǐ zǒu le liǎng ge xiǎoshí} — My granddaughter and I walked in the woods for two hours.
  \item \emph{Recall:} say the Mandarin for a tree, then the Mandarin for a fruit tree, then say \emph{shùlín} again
\end{itemize}

\begin{tcolorbox}[breakable,colback=teal!4,colframe=teal!35!black,title={Before you move on}]
What does shùlín mean? (\textquotedblleft{}Woods; a grove\textquotedblright{}.) Say it once more.
\end{tcolorbox}
\section[chéngshú]{chéngshú — to ripen; to be ripe}

\label{lesson:ZH-C516-chengshu}



\begin{quote}
Before the new one: say the Mandarin for a fruit tree, then the Mandarin for the woods.
\end{quote}

\subsection*{You'll want to know: chéngshú}
\emph{chéngshú} — \textquotedblleft{}to ripen; to be ripe\textquotedblright{}.

\begin{quote}
\emph{zhèxiē píngguǒ hái méiyǒu chéngshú} — These apples aren't ripe yet.
\end{quote}

Its characters come later; for now it is a word for the ear.

\begin{cousinweb}
\emph{chéngshú} is built from \zh{成} 'to complete' and \emph{shú}, both inherited from Middle Chinese, \emph{dzyeng} and \emph{dzyuwk} (Guangyun, in Baxter's transcription), which glosses the first 'finished' and the second with the first: done, as fruit is done when it is ripe.
\end{cousinweb}

\begin{grammarlens}[title={What you've built}]
One more word for saying how it was.
\end{grammarlens}

\subsection*{Your turn}
\begin{itemize}
\raggedright
  \item \emph{Say it:} \emph{chéngshú}
  \item \emph{Say it:} \emph{chéngshú}, once more
  \item \emph{Say it:} \emph{zhèxiē píngguǒ hái méiyǒu chéngshú} — These apples aren't ripe yet.
  \item \emph{Recall:} say the Mandarin for a fruit tree, then the Mandarin for the woods, then say \emph{chéngshú} again
\end{itemize}

\begin{tcolorbox}[breakable,colback=teal!4,colframe=teal!35!black,title={Before you move on}]
What does chéngshú mean? (\textquotedblleft{}To ripen; to be ripe\textquotedblright{}.) Say it once more.
\end{tcolorbox}
